
le composé de deux morphismes parfaits, et par suite un morphisme parfait.

\textbf{Proposition 1.10.} \uline{Soient} @@M0@@, @@M1@@ \uline{deux schémas lisses sur} @@M2@@. \uline{Alors toute} @@M3@@-\uline{immersion de} @@M4@@ \uline{dans} @@M5@@ \uline{est régulière}. \uline{En particulier}, \uline{toute section d'un morphisme lisse est une immersion régulière}.

La proposition résulte immédiatement de EGA IV 17.12.1, puisqu'une immersion régulière au sens de EGA IV est régulière au sens de 1.4.

\section*{2. Calculs sur les immersions régulières}

\subsection*{2.1.}
Soient @@M6@@ un anneau, @@M7@@ et @@M8@@ deux @@M9@@-algèbres. On se propose de munir @@M10@@ d'une structure de @@M11@@-algèbre graduée, donc de définir des homomorphismes :
@@M12@@

Soit @@M13@@ une résolution de @@M14@@ par des @@M15@@-modules projectifs; les @@M16@@ sont les groupes d'homologie du complexe @@M17@@. On sait définir des homomorphismes canoniques ([1], IV 6):
@@M18@@

Par ailleurs, @@M19@@ est un complexe de @@M20@@-modules projectifs, à degrés @@M21@@, et @@M22@@; de plus, @@M23@@ est une résolution de @@M24@@. Par suite, le morphisme d'augmentation @@M25@@ se prolonge en un morphisme de complexes : @@M26@@ et définit donc, compte tenu de l'homomorphisme @@M27@@ déduit de la structure d'algèbre de @@M28@@, des homomorphismes :
@@M29@@
